\section{Gradient descent can be perfectly stable on the data and fragile elsewhere}
\label{sec:analytic-gd-counterexample}

A global support condition need not hold for a network obtained by
ordinary gradient descent, even when the network is robust on the entire
data distribution. The following exact construction separates these
requirements without numerical attacks or a continuous-input relaxation.
It is a counterexample to an automatic guarantee, not a prevalence claim
about commonly trained networks.

\begin{proposition}[A gradient-descent trajectory with off-distribution fragility]
\label{prop:analytic-gd-counterexample}
Take $d=2$, width $N=1$, and $H\ge1$ hidden tanh layers. Let the data
distribution be uniform on $(1,1)$ and $(-1,-1)$, with label $f(x)=x_1$.
Initialize the first-layer affine map as $(x_1+x_2)/2$, all internal
scalar weights as $2$, and the readout as $v=1$. Train all weights and
the first-layer bias by full-batch gradient descent on logistic loss,
using any finite positive learning rates.

At every finite iterate, the classifier is correct on the entire data
support with signed output at least $\tanh(1)>1/2$. For every
$0<c\le1/100$, it has zero clipped-score change on this support under
every simultaneous disturbance of norm
\[
 \|\eta\|_*\le\frac{c}{\sqrt{1+Hv^2}}.
\]
The same assertion holds for the margin-band score.

At the valid Boolean input $x_0=(1,-1)$, however, its clean output is zero
and there is an admissible first-layer postactivation disturbance giving
\begin{equation}\label{eq:gd-offdata-amplification}
 z(x_0;\eta)\ge
 \min\left\{\frac12,
       \frac{c(3/2)^{H-1}}{\sqrt{H+1}}\right\}.
\end{equation}
Consequently every iterate fails the global clipped-score condition
when the second expression in the minimum exceeds $1/16$. Every iterate
fails the global margin-band condition when it exceeds $5/16$.
For $c=10^{-3}$, $H\ge16$ suffices for the former and $H\ge20$ for the
latter.
\end{proposition}

\begin{proof}
Write the first-layer weights as $(a,a)$ and its bias as zero, the
internal scalar weights as $w_s$ for $2\le s\le H$, and the readout as
$v$. At initialization $a=1/2$, $w_s=2$, and $v=1$.
The two examples have opposite inputs and labels, and tanh is odd.
Their signed outputs are equal. Their loss derivatives with respect to
the first-layer bias cancel, and the derivatives with respect to the
two first-layer weights are equal. Moreover the signed output strictly
increases with each of $a,w_2,\ldots,w_H,v$ when they are positive:
all factors in the corresponding chain derivatives are positive.
The derivative of logistic loss with respect to the signed output is
negative. Thus gradient descent preserves zero bias and equal first-layer
weights, while increasing $a,w_s,v$. In particular, throughout the
trajectory,
\[
 a\ge1/2,\qquad w_s\ge2,\qquad v\ge1.
\]
The same conclusion holds for minibatches containing equal numbers of
the two examples. It does not assume arbitrary stochastic minibatches
preserve the symmetry.

At either data input, let $y\in\{-1,1\}$ be its label. The clean first
signed hidden value is at least $\tanh(1)>3/4$. Induction gives that every subsequent signed hidden value is at least
$\tanh(1)$ because $2\tanh(1)>1$. Hence the signed readout is greater than $1/2$.

For the finite-noise assertion, each scalar disturbance, including the
output disturbance, has magnitude at most $c$. The first noisy signed
hidden state is at least $\tanh(1)-c>0.7$. If the preceding noisy signed
state is greater than $0.7$, then the next is at least
$\tanh(1.4)-c>\tanh(1)-c>0.7$.
Induction gives a final signed output greater than
$0.7v-c\ge0.69>1/2$. Clean and noisy outputs therefore lie beyond the
same clipping endpoint. Both the clipped score and margin-band score
are exactly unchanged on the data support.

At $x_0$, the equal first-layer weights cancel and every clean hidden
state is zero. Inject only at the first hidden layer, with
\[
 \eta_1=\frac{c}{v\sqrt{H+1}}.
\]
This is admissible because
$(1+Hv^2)\eta_1^2\le c^2$ follows from $v\ge1$.
All subsequent noisy hidden states are nonnegative. Concavity of tanh
on the positive half-line and $\tanh(1)>3/4$ imply
\[
 \tanh(2u)\ge\min\{1/2,(3/2)u\}\qquad(u\ge0).
\]
For $0\le u\le1/2$, use the chord bound
$\tanh(2u)\ge2u\tanh(1)\ge(3/2)u$; for larger $u$, use
$\tanh(2u)\ge\tanh(1)>1/2$.
Since $w_s\ge2$, iterating the displayed bound yields
\[
 z(x_0;\eta)\ge
 \min\left\{v/2,\frac{c(3/2)^{H-1}}{\sqrt{H+1}}\right\},
\]
which proves~\eqref{eq:gd-offdata-amplification}.
The clean transformed scores are zero. A noisy output greater than
$1/16$ violates the clipped-score tolerance; an output greater than
$1/4+1/16=5/16$ violates the margin-band tolerance.
At $c=10^{-3}$, the latter expression is greater than $0.106$ at
$H=16$ and greater than $0.483$ at $H=20$. It increases with $H$,
proving the stated sufficient depths.
\end{proof}

This example also shows why the trusted-region variant is substantive.
Taking the trusted region to be the two-point data support gives an exact
whole-region stability certificate for every gradient-descent iterate,
whereas a full-cube condition rejects all these same networks at sufficient
depth. Extending this observation to realistic data requires an actual
region-wide certificate or a controlled family of exceptional regions;
held-out pointwise success alone is not such a certificate.
